%! AP Statistics — Unit 2, Lesson 2.2 — Exit Ticket KEY
\documentclass[10pt]{article}
\usepackage{apstats-article}
\usepackage{apstats-key}

%=============================================================================
\begin{document}

\pageheader{Unit 2, Lesson 2.2}{Exit Ticket --- Answer Key}
\begin{tabularx}{\linewidth}{X X X}
Name:\;\ans{Key} & Date:\; & Period:\;
\end{tabularx}

\vspace{0.22in}

\begin{tcolorbox}[colback=sky, colframe=navy, boxrule=0.9pt, arc=2mm,
    left=4mm, right=4mm, top=3mm, bottom=3mm]
\small\textit{A researcher surveyed 180 students about whether they own a pet (Yes/No)
and whether they volunteer in their community (Yes/No).}

\vspace{10pt}
\renewcommand{\arraystretch}{1.9}
\begin{tabularx}{\linewidth}{@{} l C{2.6cm} C{2.6cm} C{2.4cm} @{}}
\rowcolor{sky}
                    & \textbf{Volunteer: Yes} & \textbf{Volunteer: No} & \textbf{Row Total} \\
\hline
Pet: Yes   & 48 & 72 & \ans{120} \\
Pet: No    & 24 & 36 & \ans{60} \\
\hline
Col.\ Total& \ans{72} & \ans{108} & \ans{180} \\
\end{tabularx}
\end{tcolorbox}

\vspace{0.18in}

\begin{enumerate}[leftmargin=1.6em, itemsep=0.18in]

    \item \ans{Row totals: 120, 60. Column totals: 72, 108. Grand total: 180.}

    \item P(Vol $|$ Pet Yes) $= \dfrac{\ans{48}}{\ans{120}} = \ans{0.40 = 40\%}$
          \hfill
          P(Vol $|$ Pet No) $= \dfrac{\ans{24}}{\ans{60}} = \ans{0.40 = 40\%}$

    \item Circle: \textbf{No} \textcolor{keyred}{\textbf{$\leftarrow$ correct}}\\[8pt]
          \ans{The conditional relative frequency of volunteering is 40\% for
          \emph{both} pet owners and non-pet owners. Because the conditional
          distributions are identical, there is \textbf{no} evidence of an
          association between pet ownership and volunteering.}

\end{enumerate}

\vspace{0.2in}

\begin{tcolorbox}[colback=redbg, colframe=redacc, boxrule=0.8pt, arc=2mm,
    left=4mm, right=4mm, top=3mm, bottom=3mm]
\small\textbf{AP-Style Practice}\\[6pt]
9th: 60\% participate; 10th: 55\% participate.

\begin{enumerate}[label=\textbf{(\Alph*)}, leftmargin=2em, itemsep=6pt]
  \item There is no association because both values are below 100\%.
  \item There is a strong positive association because more than half participate.
  \item There is weak or no association because the conditional distributions are nearly the same.
        \textcolor{keyred}{\textbf{$\leftarrow$ correct}}
  \item Association cannot be determined without the actual cell counts.
\end{enumerate}

\vspace{6pt}
\ans{(C). A difference of 5 percentage points (60\% vs.\ 55\%) is very small --- the
conditional distributions are nearly identical, so there is little or no evidence of
an association. (D) is a common distractor --- we \emph{can} use conditional relative
frequencies, which are given, to judge association.)
}
\end{tcolorbox}

\begin{teachernote}
\textbf{Item 3 (No association):} Many students expect the answer to be ``Yes'' and will be
surprised that the conditional distributions are identical. Use this as a teaching moment:
identical conditional distributions $=$ no association, regardless of the raw counts.
\textbf{AP item:} (B) is the most common wrong answer --- students confuse the direction of
an association with whether one exists. Emphasize that ``association'' means the conditional
distributions \emph{differ}, not simply that participation is high.
\end{teachernote}

\end{document}
